% Generated draft: requires mathematical and source review.
\documentclass[11pt]{article}
\usepackage[margin=25mm]{geometry}
\usepackage{fontspec}
\setmainfont{Latin Modern Roman}
\usepackage{amsmath,amssymb,mathtools,unicode-math}
\setmathfont{Latin Modern Math}
\usepackage{enumitem,needspace,xurl,hyperref}
\setlength{\parindent}{0pt}
\setlength{\parskip}{5pt}
\setlength{\emergencystretch}{3em}
\newcommand{\N}{\mathbb{N}}
\newcommand{\Z}{\mathbb{Z}}
\newcommand{\Q}{\mathbb{Q}}
\newcommand{\R}{\mathbb{R}}
\newcommand{\C}{\mathbb{C}}
\newcommand{\F}{\mathbb{F}}
\newcommand{\A}{\mathbb{A}}
\newcommand{\PP}{\mathbb{P}}
\newcommand{\E}{\mathbb{E}}
\newcommand{\eps}{\varepsilon}
\newcommand{\dd}{\,\mathrm{d}}
\begin{document}
\section*{UnsolvedMath ID: 30001594}
\textbf{Source problem code:} catalogue entry
\subsection{Definitions and mathematical statement}
The supplied record asks: \emph{What are caustics in the oscillatory-integral setting, and how can their geometry and order be determined?} It describes caustics as envelopes of optical rays or as degeneracies of stationary points of a phase, associated with failure of uniform \(L^\infty\) bounds. This is a research question whose mathematical setting is not specified sufficiently to give a unique quantified statement.

To explain the phase terminology only, let \(X\subseteq\R^d\) and \(\Theta\subseteq\R^m\) be open sets, with \(d,m\geq1\), and let \(\phi\in C^2(X\times\Theta;\R)\). These are auxiliary conventions, not hypotheses supplied for the intended problem. A stationary point in the integration variables is a pair \((x,\theta)\) satisfying
\[
\nabla_\theta\phi(x,\theta)=0.
\]
It is degenerate when
\[
\det D^2_{\theta\theta}\phi(x,\theta)=0.
\]
The corresponding set of parameter values is
\[
\mathcal K_\phi=\bigl\{x\in X:\ \exists\theta\in\Theta,\ \nabla_\theta\phi(x,\theta)=0,\ \det D^2_{\theta\theta}\phi(x,\theta)=0\bigr\}.
\]
This set expresses a possible local phase description of a caustic. Identifying it with an envelope of optical rays requires additional geometric hypotheses and a specified ray family, neither of which is supplied.

\textbf{Statement of the supplied research target.} Determine the geometry and order of caustics in the intended oscillatory-integral or optical-ray setting, and characterize the associated failure of the intended uniform \(L^\infty\) estimates. The source does not specify the family of integrals, phase class, amplitude, parameter regime, ray dynamics, domains, boundary conditions, or meaning of ``order.'' It also does not identify which bound fails or what its uniformity means. Degeneracy alone does not imply failure of every uniform \(L^\infty\) bound; the relevant normalization, amplitude, and limiting regime must be specified. Thus the displayed definitions clarify terminology but do not complete the missing hypotheses or establish a classification theorem.
\subsection{Short English statement}
The question asks for the geometry and order of caustics, described through optical-ray envelopes or degenerate stationary points of oscillatory phases. The supplied fragment leaves the class of examples, the meaning of order, and the intended uniform \(L^\infty\) estimate unspecified, so a precise classification problem cannot yet be extracted.
\subsection{Sources}
\begin{itemize}
\item Supplied local UnsolvedMath registry record: permanent ID \(30001594\), title ``Geometry and Order of Caustics,'' source problem code \texttt{catalogue entry}, in the Oberwolfach Reports Open Problems collection. The proposer is not identified; the supplied proposed year is 2010.
\item Supplied source citation: \emph{Nonlinear Waves and Dispersive Equations} (2011), Oberwolfach Reports, DOI \texttt{10.4171/owr/2010/41}, pp.\ 2393--2463. \url{https://doi.org/10.4171/owr/2010/41}. The collection background describes Oberwolfach Reports as documenting mathematical workshops and their open-problem discussions; no particular page or speaker is identified for this fragment.
\item Additional supplied literature reference: \url{https://www.numdam.org/article/JEDP_2010____A11_0.pdf}. No author, title, or pinpoint reference is supplied for this document.
\item Literature assessment dated exactly \texttt{2026-08-21}: the supplied curation labels the problem ``open'' and characterizes the extraction as a fragment requesting a general theory of caustics. It states that existing results classify important canonical singularities but do not turn this fragment into a single solved problem. This is supplied dated curation, not an independently verified literature assessment; the record explicitly calls for independent verification of time-sensitive status and cited claims.
\item The registry labels the statement \texttt{corrected\_verified} and reports a manual semantic, typographic, and LaTeX review claiming a self-contained, meaning-preserving cleanup without unsupported additions. These are supplied editorial assertions, not independent verification of the formulation or its mathematical status.
\end{itemize}
\end{document}
